\chapter[1982 --- Bergstrom et al.]{1982: Sune K. Bergstr\"om, Bengt I. Samuelsson, John R. Vane}
\label{ch:1982_Bergstrom}

\section*{Main Contribution}

\textit{Discovered prostaglandins and related biologically active substances, revealing a new class of chemical messengers that regulate inflammation, blood flow, and many other functions.}

\section{Official Citation}

for their discoveries concerning prostaglandins and related biologically active substances

The 1982 Nobel Prize in Physiology or Medicine was awarded jointly to Sune K. Bergstr\"om, Bengt I. Samuelsson, and John R. Vane for their discoveries concerning prostaglandins and related biologically active substances. Their work transformed our understanding of lipid signaling molecules and led to the development of some of the most widely used drugs in modern medicine, including aspirin-like anti-inflammatory agents, prostaglandin analogs for the treatment of cardiovascular and reproductive disorders, and drugs for the management of asthma and other inflammatory conditions.

\section{Background and Historical Context}

Prostaglandins were first described in the 1930s by the Swedish physiologist Ulf von Euler, who observed that seminal fluid contained a substance that caused smooth muscle contraction and lowered blood pressure. Von Euler named the substance ``prostaglandin'' because he initially believed it was produced by the prostate gland (in fact, prostaglandins are produced by many tissues throughout the body). However, the chemical structure of prostaglandins remained unknown for decades, and the full extent of their biological activities was not appreciated until the work of Bergstr\"om, Samuelsson, and Vane.

The study of prostaglandins was complicated by their notable potency and the vanishingly small quantities in which they are produced. Prostaglandins are active at nanomolar or even picomolar concentrations, and they are rapidly metabolized and inactivated by the body. Isolating and characterizing these substances required the development of new biochemical techniques and the processing of enormous quantities of biological materials.

Sune Karl Bergstr\"om was born in Stockholm, Sweden, in 1916, and trained in medicine and biochemistry at the Karolinska Institute. Bergstr\"om became interested in the chemical nature of prostaglandins in the 1950s, and he devoted much of his career to isolating and characterizing these substances.

Bengt Ingemar Samuelsson was born in Halmstad, Sweden, in 1934, and trained in organic chemistry and biochemistry at the University of Lund. He joined Bergstr\"om's laboratory at the Karolinska Institute in the early 1960s and became one of the leading figures in prostaglandin research.

John Robert Vane was born in Tardebigg, Worcestershire, England, in 1927, and trained in chemistry at the University of Birmingham. He joined the Wellcome Foundation in London in 1953 and later moved to the Institute of Basic Medical Sciences at the University of London, where he conducted his Nobel Prize-winning work on prostaglandins and aspirin.

\section{The Discovery/Contribution}

\subsection{Sune K. Bergstr\"om and the Chemical Characterization of Prostaglandins}

Bergstr\"om's contribution was the chemical characterization of the major prostaglandins. Using a combination of chemical extraction, chromatography, and mass spectrometry, Bergstr\"om and his colleagues isolated and purified prostaglandins from sheep seminal vesicles and other tissues. They determined the chemical structures of the major prostaglandins---PGE$_1$, PGF$_{1\alpha}$, PGE$_2$, and PGF$_{2\alpha}$---and showed that all prostaglandins share a common basic structure: a five-membered ring (cyclopentane ring) with two side chains. The prostaglandins differ in the pattern of hydroxyl and keto groups on the ring and in the degree of unsaturation of the side chains.

Bergstr\"om's structural work was a tour de force of organic chemistry. The prostaglandins are complex molecules with multiple stereocenters, and determining their structures required the application of advanced analytical techniques, including infrared spectroscopy, nuclear magnetic resonance spectroscopy, and mass spectrometry. Bergstr\"om also synthesized several prostaglandin analogs, which were used to study the structure-activity relationships of the prostaglandins and to develop drugs based on the prostaglandin structure.

Bergstr\"om's work established the biochemical foundation for all subsequent prostaglandin research. By determining the chemical structures of the prostaglandins, Bergstr\"om made it possible to synthesize these molecules in the laboratory, to develop assays for their measurement in biological fluids, and to design drugs that mimic or block their actions.

\subsection{Bengt I. Samuelsson and the Biosynthesis of Prostaglandins}

Samuelsson's contribution was the elucidation of the biosynthetic pathway for prostaglandins. Working in Bergstr\"om's laboratory and later at the Karolinska Institute, Samuelsson discovered that prostaglandins are synthesized from a common precursor---arachidonic acid, a polyunsaturated fatty acid that is released from cell membrane phospholipids by the enzyme phospholipase A$_2$ in response to various stimuli (including tissue injury, inflammation, and hormonal signals).

Samuelsson showed that arachidonic acid is converted to a labile intermediate called prostaglandin endoperoxide (PGG$_2$ and PGH$_2$) by the enzyme cyclooxygenase (COX). The endoperoxides are then converted to the various prostaglandins and thromboxanes by specific isomerases and reductases. Samuelsson also discovered that the prostaglandin endoperoxides are precursors of another important class of biologically active substances---the thromboxanes, which are produced by platelets and play a crucial role in blood clotting and platelet aggregation.

Perhaps Samuelsson's primary discovery was the identification of the leukotrienes---a family of potent inflammatory mediators that are produced from arachidonic acid by a pathway parallel to the prostaglandin pathway. Samuelsson showed that leukotrienes are produced by the enzyme 5-lipoxygenase and that they are responsible for the bronchoconstriction and mucus secretion that occur in asthma. This discovery had immediate clinical implications: it led to the development of leukotriene receptor antagonists (such as montelukast/Singulair) and leukotriene synthesis inhibitors (such as zileuton) for the treatment of asthma.

\subsection{John R. Vane and the Mechanism of Action of Aspirin}

Vane's contribution was the discovery of the mechanism by which aspirin and other non-steroidal anti-inflammatory drugs (NSAIDs) exert their therapeutic effects. In 1971, Vane published a landmark paper in \emph{Nature} showing that aspirin and indomethacin inhibit the biosynthesis of prostaglandins by blocking the activity of cyclooxygenase (COX). This was a discovery of immense practical importance: it explained how aspirin---one of the most widely used drugs in the world---works at the molecular level.

Vane's discovery was made using a bioassay system in which the release of prostaglandins from guinea pig lungs was measured by its effect on smooth muscle preparations. He showed that aspirin blocked the release of prostaglandins from the lungs in a dose-dependent manner, and he proposed that aspirin acted by inhibiting the enzyme that converts arachidonic acid to prostaglandin endoperoxides. This enzyme was later purified and named cyclooxygenase (COX).

Vane's discovery of the mechanism of action of aspirin had several important consequences. First, it provided a rational basis for the use of aspirin as an anti-inflammatory, analgesic, and antipyretic agent. Second, it led to the development of new NSAIDs (such as ibuprofen, naproxen, and diclofenac) that act by the same mechanism---inhibition of COX---but with different pharmacological profiles. Third, it provided the basis for the development of selective COX-2 inhibitors (such as celecoxib and rofecoxib), which were designed to inhibit the inducible form of COX (COX-2) that is responsible for inflammation, while sparing the constitutive form (COX-1) that protects the stomach lining from acid. (The COX-2 inhibitor rofecoxib/Vioxx was later withdrawn from the market due to cardiovascular side effects, but celecoxib/Celebrex remains in use.)

Vane's work also contributed to the understanding of the role of prostaglandins in the cardiovascular system. He discovered that the endothelium (the inner lining of blood vessels) produces a substance that relaxes vascular smooth muscle---a substance that was later identified as nitric oxide (NO) by Robert Furchgott, Louis Ignarro, and Ferid Murad (who shared the 1998 Nobel Prize in Physiology or Medicine for this discovery). Vane's work on endothelial-derived relaxing factor (EDRF) and nitric oxide complemented his prostaglandin research and further elucidated the mechanisms by which the cardiovascular system maintains normal blood flow and prevents clotting.

\section{Impact on Modern Medicine}

The discoveries of Bergstr\"om, Samuelsson, and Vane have had an enormous impact on modern medicine. The most immediate practical consequence was the elucidation of the mechanism of action of aspirin, which has become one of the most widely used drugs in the world. Aspirin is used as an analgesic, antipyretic, and anti-inflammatory agent, and it is also used in low doses for the prevention of cardiovascular events (heart attacks and strokes). The cardiovascular benefits of low-dose aspirin are a direct consequence of its ability to inhibit thromboxane synthesis in platelets, thereby reducing platelet aggregation and the risk of thrombosis.

The development of NSAIDs---including ibuprofen, naproxen, and diclofenac---based on the understanding of COX inhibition has provided millions of patients with effective relief from pain, inflammation, and fever. These drugs are among the most commonly prescribed medications worldwide and are used in the management of conditions ranging from arthritis to postoperative pain to menstrual cramps.

Samuelsson's discovery of the leukotrienes led to the development of leukotriene receptor antagonists (such as montelukast) for the treatment of asthma. These drugs provide an alternative to inhaled corticosteroids for the management of mild to moderate asthma, and they have improved the quality of life of millions of asthma patients worldwide.

Prostaglandin analogs have also found important clinical applications. Prostaglandin E$_2$ (dinoprostone) is used to induce labor and to ripen the cervix before delivery. Prostacyclin (PGI$_2$) analogs (such as epoprostenol and iloprost) are used to treat pulmonary hypertension, a potentially fatal condition in which high blood pressure in the lungs leads to right heart failure. Misoprostol, a prostaglandin E$_1$ analog, is used to prevent gastric ulcers in patients taking NSAIDs, and it is also used for the medical management of miscarriage and for the prevention of postpartum hemorrhage.

The elucidation of the arachidonic acid cascade---the complex network of enzymes and signaling molecules that metabolize arachidonic acid---has also provided the basis for the development of new drugs for the treatment of inflammatory diseases, cardiovascular disease, and cancer. The recognition that COX-2 is overexpressed in many cancers has led to the investigation of COX-2 inhibitors as potential anticancer agents, and there is growing evidence that aspirin and other NSAIDs may reduce the risk of colorectal cancer and other malignancies.

\section{Legacy}

Sune Bergstr\"om served as rector of the Karolinska Institute and was instrumental in the development of prostaglandin research as a major field of biochemistry. He died in 2004, but his chemical characterization of the prostaglandins continues to serve as the foundation of prostaglandin research and drug development.

Bengt Samuelsson continued his work at the Karolinska Institute, where he served as rector and made further contributions to our understanding of leukotrienes, lipoxins, and other lipid mediators. He has been a prominent advocate for international scientific cooperation and for the application of basic research to clinical medicine.

John Vane continued his work at the Wellcome Research Laboratories, where he directed research on prostaglandins, nitric oxide, and other signaling molecules. He died in 2004, but his discovery of the mechanism of action of aspirin---one of the major drug mechanisms ever elucidated---continues to influence the practice of medicine.

The work of the three laureates of 1982 demonstrated the power of biochemical and pharmacological approaches to understand disease and to develop new therapies. Their discoveries about prostaglandins and related substances have saved and improved countless lives, and they continue to inspire new research into the complex signaling networks that regulate inflammation, cardiovascular function, and reproduction.


\section{Modern Developments}

Bergström, Samuelsson, and Vane's prostaglandin work has led to modern anti-inflammatory therapy. COX-2 inhibitors (celecoxib) reduce inflammation without gastrointestinal side effects. Understanding of prostaglandins has led to treatments for pulmonary hypertension (epoprostenol), glaucoma (latanoprost), and erectile dysfunction (alprostadil). Understanding of thromboxane has led to antiplatelet therapies (aspirin, clopidogrel) for cardiovascular disease prevention.


\section{Bibliography}

\begin{thebibliography}{9}
\bibitem{nobel-1982}
Nobel Prize Outreach, ``The Nobel Prize in Physiology or Medicine 1982,''\emph{NobelPrize.org}.\\
\url{https://www.nobelprize.org/prizes/medicine/1982/summary/}

\bibitem{lecture-1982}
Sune K. Bergstr"om, ``Nobel Lecture,''\emph{NobelPrize.org}. Winner-authored primary source; downloadable from the official Nobel Prize archive.\\
\url{https://www.nobelprize.org/prizes/medicine/1982/bergstrom/lecture/}
\end{thebibliography}
